% Part: first-order-logic
% Chapter: introduction
% Section: soundness-completeness

\documentclass[../../../include/open-logic-section]{subfiles}

\begin{document}

\olfileid{fol}{int}{scp}

\olsection{Correctheid en volledigheid}

We voeren ook !!{derivation}ssystemen voor de eerste-ordelogica in.
Logici hebben veel verschillende !!{derivation}ssystemen bedacht. Toch
geven ze allemaal dezelfde !!{derivability}srelatie tussen
!!{sentence}s. We zeggen dat uit~$\Gamma$ de formule~$!A$ valt te
\emph{!!{derive}}, en schrijven $\Gamma \Proves !A$, als er een
!!a{derivation} van een precies omschreven soort bestaat.
!!^{derivation}s zijn altijd eindige patronen van symbolen. Zo'n
patroon kan een lijst !!{sentence}s zijn, maar ook iets ingewikkelders.
Een !!{derivation}ssysteem moet ons helpen bepalen of !!a{sentence} uit
een verzameling~$\Gamma$ volgt. Daarom moet gelden:
$\Gamma \Entails !A$ precies als $\Gamma \Proves !A$.

Stel dat $\Gamma \Proves !A$ wel geldt, maar $\Gamma \Entails !A$
niet. Dan is ons !!{derivation}ssysteem te sterk: het bewijst te veel.
Een systeem is \emph{correct} als uit $\Gamma \Proves !A$ altijd
$\Gamma \Entails !A$ volgt. Dat is de minste eis waaraan een goed
!!{derivation}ssysteem moet voldoen. Stel nu dat $\Gamma \Entails !A$
wel geldt, maar $\Gamma \Proves !A$ niet. Dan is het systeem te zwak:
het bewijst te weinig. Een systeem is \emph{volledig} als uit
$\Gamma \Entails !A$ altijd $\Gamma \Proves !A$ volgt. Correctheid is
meestal vrij eenvoudig te bewijzen. We gebruiken inductie naar de
opbouw van de inductief gedefinieerde !!{derivation}s. Volledigheid is
moeilijker.

Correctheid en volledigheid hebben belangrijke gevolgen. Als uit een
verzameling !!{sentence}s~$\Gamma$ een tegenspraak, zoals
$!A \land \lnot !A$, valt te !!{derive}, heet~$\Gamma$
\emph{inconsistent}. Een inconsistente verzameling kan geen model
hebben en is dus onvervulbaar. Uit volledigheid volgt ook het omgekeerde.
Iedere~$\Gamma$ die niet inconsistent is---we noemen haar dan
\emph{consistent}---heeft een model. Dit is zelfs een andere,
gelijkwaardige vorm van volledigheid. Die vorm gaan we bewijzen.
Het resultaat is diep en misschien verrassend. Kun je de formule
$!A \land \lnot !A$ niet uit~$\Gamma$ bewijzen, dan moet er
!!a{structure} bestaan die is zoals~$\Gamma$ haar beschrijft. De vraag
welke verzamelingen !!{sentence}s modellen hebben, heeft dus een kort
antwoord: precies de consistente verzamelingen.

De correctheids- en volledigheidsstelling hebben op hun beurt twee
belangrijke gevolgen: de compactheidsstelling en de stelling van
L\"owenheim--Skolem. Met deze resultaten uit de modeltheorie kunnen we
weer veel andere zaken bewijzen. Twee voorbeelden kwamen al voorbij.
De eerste-ordelogica kan niet zeggen dat het !!{domain} van
!!a{structure} eindig is. Zij kan ook niet zeggen dat het
!!{nonenumerable} is.

Historisch kostte het veel tijd om dit allemaal goed uit te werken:
de syntaxis en semantiek van de eerste-ordelogica, goede
!!{derivation}ssystemen, bewijzen van correctheid en volledigheid en de
grens tussen wat eerste-ordetalen wel en niet kunnen uitdrukken. Nu
weten we hoe het moet. Toch kan het nog steeds verwarrend en taai zijn
om alle details door te nemen. Die details zijn wel belangrijk. De
methoden voor de formele taal van de eerste-ordelogica worden overal
gebruikt: in de logica, informatica en taalkunde. Wie ze nu doorwerkt,
heeft daar later voordeel van.

\end{document}
